% =====================================================================
% packages.tex - Tất cả các package cần thiết
% =====================================================================

% ---------------------- NGÔN NGỮ & MÃ HÓA ----------------------------
% Hỗ trợ đúng phông tiếng Việt với cả pdfLaTeX và XeLaTeX.
\usepackage{iftex}
\ifPDFTeX
    \usepackage[utf8]{inputenc}
    \usepackage[T5]{fontenc}
    \usepackage[utf8]{vietnam}
\else
    \usepackage{fontspec}
    \defaultfontfeatures{Ligatures=TeX}
    \setmainfont{Times New Roman}
\fi

% ---------------------- BỐ CỤC TRANG ---------------------------------
\usepackage[top=2.5cm, bottom=2.5cm, left=3.5cm, right=2.5cm]{geometry}
\usepackage{fancyhdr}                % Tùy chỉnh header/footer
\usepackage{setspace}                % Giãn dòng

% ---------------------- HIỂN THỊ MÃ NGUỒN ---------------------------
\usepackage{listings}                % Hiển thị code
\usepackage{xcolor}                  % Màu sắc

% ---------------------- BIỂU ĐỒ & HÌNH VẼ --------------------------
\usepackage{tikz}                    % Vẽ hình
\usepackage{graphicx}                % Chèn ảnh ERD

% ---------------------- BẢNG BIỂU --------------------------------
\usepackage{booktabs}                % Bảng chuyên nghiệp
\usepackage{longtable}               % Bảng dài (vượt nhiều trang)
\usepackage{array}                   % Tùy chỉnh căn lề và định dạng cột

% ---------------------- DANH SÁCH & LIỆT KÊ ----------------------------
\usepackage{enumitem}                % Tùy chỉnh danh sách

% ---------------------- TIÊU ĐỀ & PHẦN ----------------------------
\usepackage{titlesec}                % Tùy chỉnh tiêu đề chapter/section
\usepackage{tocloft}                 % Tùy chỉnh cỡ chữ và thụt lề mục lục

% ---------------------- LIÊN KẾT ----------------------------
\usepackage{hyperref}                % Tạo liên kết

% =====================================================================
% TÙY CHỈNH THƯ VIỆN
% =====================================================================
\usetikzlibrary{positioning, shapes, arrows, calc}
